\chapter*{پیشگفتار}
\phantomsection
\addcontentsline{toc}{chapter}{پیشگفتار}
\markboth{پیشگفتار}{پیشگفتار}

بینایی ماشین یکی از حوزه‌های محوری هوش مصنوعی است که هدف آن استخراج اطلاعات معنادار از تصویر و ویدئو است. آشکارسازی شیء، علاوه بر تعیین رده اشیا، مکان آن‌ها را نیز مشخص می‌کند و به همین دلیل در سامانه‌های نظارت هوشمند، رباتیک، حمل‌ونقل خودکار، بازرسی صنعتی و تحلیل ویدئو کاربرد دارد. یک آشکارساز عملی باید در برابر تغییر مقیاس، انسداد، تراکم اشیا، تغییر نور و پس‌زمینه‌های پیچیده پایدار باشد و هم‌زمان با تأخیر محدود اجرا شود \cite{lin2014coco,ren2015faster}.

آشکارسازهای تک‌مرحله‌ای با پیش‌بینی مستقیم جعبه‌های مرزی و برچسب‌های رده از نقشه‌های ویژگی، مسیر مؤثری برای دستیابی به سرعت بالا ایجاد کرده‌اند. خانواده YOLO نمونه برجسته این رویکرد است و در نسخه‌های گوناگون، طراحی ستون فقرات، گردن و سر آشکارساز را برای بهبود توازن دقت و سرعت تکامل داده است \cite{redmon2016yolo,liu2016ssd}. YOLOv13 نیز با استفاده از کانولوشن‌های تفکیک‌پذیر و سازوکارهای تقویت و توزیع اطلاعات چندمقیاسی، یک خط پایه مناسب برای پژوهش حاضر فراهم می‌کند \cite{lei2025yolov13}.

در بسیاری از صحنه‌ها، تشخیص قابل اعتماد به دریافت هم‌زمان جزئیات محلی و ارتباط‌های دوربرد نیاز دارد. مبدل‌ها با سازوکار توجه، در مدل‌سازی وابستگی‌های دوربرد توانمند هستند؛ اما توجه سراسری در وضوح بالا هزینه محاسبات را افزایش می‌دهد \cite{vaswani2017attention,dosovitskiy2021vit}. معماری «مبدل در مبدل» (\lr{TNT}) با تفکیک تعاملات درون‌قطعه‌ای و میان‌قطعه‌ای، راهی برای سازمان‌دهی توجه محلی و سراسری ارائه می‌دهد \cite{han2021tnt}. از سوی دیگر، بلوک‌های Inception و توجه کانالی کارا با ترکیب میدان‌های دید متفاوت و بازوزن‌دهی کانال‌ها، ابزارهایی مناسب برای استخراج ویژگی چندمقیاسی با سربار محدود هستند \cite{szegedy2015inception,wang2020eca}.

این پایان‌نامه بر طراحی یک آشکارساز سبک و دقیق مبتنی بر YOLOv13 متمرکز است. در این مسیر، نقش توجه سلسله‌مراتبی، بلوک‌های چندشاخه و توجه هندسی در کیفیت بازنمایی و هزینه اجرا بررسی می‌شود. فصل اول مسئله، انگیزه، اهداف، پرسش‌ها و چارچوب ارزیابی پژوهش را بیان می‌کند. فصل‌های بعدی به مبانی نظری، معماری، آزمایش‌ها و جمع‌بندی اختصاص خواهند داشت.
